% Source: https://artofproblemsolving.com/community/c3h1657695_2018_ms_mte_theme_round_all
% Stable identity: aops-topic-1657695-q06
By the shoelace formula the area is
\[
\dfrac{1}{2}\left|x(3x-5x)+18(5x-2)+6(2-3x)\right|=\dfrac{1}{2}\left|-2x^2+72x-24\right|=\left|-x^2+36x-12\right| .
\]
The polynomials whose absolute value equals this for all $x$ are $\pm\left(-x^2+36x-12\right)$; the one with $p(0)<0$ is $p(x)=-x^2+36x-12$. The sum of its coefficients is $p(1)=-1+36-12=\boxed{23}$. For $x=1$ the triangle is the one shown, with area $23$.

\begin{asy}
usepackage("amsmath");
size(7cm);
defaultpen(fontsize(10pt));
pair P1=(1,2),P2=(18,3),P3=(6,5);
draw((-0.5,0)--(19.5,0),Arrow(6));draw((0,-0.5)--(0,6.5),Arrow(6));
label("$x$",(19.5,0),SE);label("$y$",(0,6.5),NW);
draw(P1--P2--P3--cycle,linewidth(0.9));
dot(P1);dot(P2);dot(P3);
label("$(1,2)$",P1,SE);label("$(18,3)$",P2,SE);label("$(6,5)$",P3,N);
label("\mbox{}",(min(currentpicture,user=true)+max(currentpicture,user=true))/2,black); // keeps the text after the figure black when LaTeX typesets the labels inline
\end{asy}
